\documentclass[12pt]{article}
\usepackage[utf8]{inputenc}
\usepackage[T1]{fontenc}
\usepackage{hyperref}
\title{Tii Language Resources for K-12 Curriculum Development}
\author{Compiled by Claude Sonnet 5 (claude-sonnet-5) for Andr\'as Kornai}
\date{}
\begin{document}
\maketitle
\section{Introduction}
Tii (ISO 639-3 txq, Glottocode tiii1241) is an Austronesian language of Rote Island, Indonesia, in the Rote-Meto subgroup. Wikipedia gives about 20,000 speakers (2002). Searches found very little material specific to Tii.
\section{Primary Educational Resources}
None found in search results.
\section{Linguistic and Reference Materials}
\begin{itemize}
\item Wikipedia, Tii language (a stub): \url{https://en.wikipedia.org/wiki/Tii_language}
\item Edwards, Rote-Meto Comparative Dictionary (ANU Press, 2021; a reconstruction of the Rote-Meto group, not a Tii dictionary): \url{https://library.oapen.org/handle/20.500.12657/52033}
\item Endangered Languages Project bibliography: \url{https://www.endangeredlanguages.com/lang/9129/bibliography?hl=de}
\end{itemize}
\section{Religious and Community Resources}
Joshua Project reports that Tii speakers are mostly Christian and have no Bible in their own language, using the Indonesian Bible instead: \url{https://joshuaproject.net/people_groups/19494}
\section{Digital and Audio Resources}
None found in search results.
\section{Curriculum Development Context}
No orthography, school programme or native-language teaching material was found. Development would start from new documentation.
\end{document}